% !TEX root = ../article.tex
% =============================================================================
% Preamble - Clase, paquetes y formato del informe
% Thoraxis - Proyecto Integrador del Programa en Deep Learning (UTEC Posgrado)
% =============================================================================
% Artículo a una sola columna, sin límite de páginas. La guía del proyecto
% integrador (Documentos/Proyecto-Integrador-Grupal-Deep-Learning.pdf) fija
% las secciones y que solo el resumen ejecutivo no exceda una página, pero no
% el formato visual.
% El aspecto (Century Schoolbook) se toma de la tesis del
% autor (github.com/ahincho/thesis) para mantener un estilo propio y
% coherente.
%
% Todo lo que define el aspecto vive en este archivo; el texto vive en
% sections/.
% =============================================================================

\documentclass[11pt, a4paper, oneside, onecolumn]{article}

% --- Codificación ---
% Los .tex se escriben en UTF-8 con acentos directos, sin \'a.
\usepackage[utf8]{inputenc}
\usepackage[T1]{fontenc}

% --- Idioma ---
%   es-noquoting     desactiva el " activo de babel-spanish (choca con tablas)
%   es-tabla         "Tabla" en lugar de "Cuadro"
%   es-nodecimaldot  no cambia el punto decimal en modo matemático
\usepackage[english, spanish, es-noquoting, es-tabla, es-nodecimaldot]{babel}

% --- Geometría de página ---
% headheight da espacio al logo del encabezado (\articleHeaderLogoHeight).
\usepackage[top=2.5cm, bottom=2.5cm, left=2.5cm, right=2.5cm, headheight=22pt]{geometry}

% --- Tipografía ---
% TeX Gyre Schola (Century Schoolbook) con matemáticas a juego.
\usepackage{tgschola}
\usepackage[nc]{newtxmath}
\usepackage{microtype}

% --- Interlineado y párrafos ---
\usepackage{setspace}
\setstretch{1.15}
\setlength{\parindent}{1.25em}
\setlength{\parskip}{0.3em}

% Válvula de escape contra Overfull \hbox en términos técnicos largos.
\setlength{\emergencystretch}{1.5em}
\hyphenation{PyTorch MONAI FastAPI NiiVue Cornerstone Evidently GitHub
    ReproPath Thoraxis}

% --- Colores ---
\usepackage{xcolor}
\definecolor{articlegray}{gray}{0.35}   % Títulos de sección
\definecolor{articlelink}{HTML}{0B5394} % Enlaces

% --- Gráficos ---
\usepackage{graphicx}
\graphicspath{{figures/}}
% Los diagramas se hacen con Excalidraw y se insertan como PDF (ver diagrams/).

% --- Títulos de secciones ---
\usepackage{titlesec}
\titleformat{\section}{\normalfont\Large\bfseries\color{articlegray}}{\thesection.}{0.6em}{}
\titleformat{\subsection}{\normalfont\large\bfseries}{\thesubsection.}{0.6em}{}
\titleformat{\subsubsection}{\normalfont\normalsize\bfseries}{\thesubsubsection.}{0.6em}{}
\titlespacing*{\section}{0pt}{1.4em}{0.8em}
\titlespacing*{\subsection}{0pt}{1.1em}{0.5em}
\setcounter{secnumdepth}{3}
\setcounter{tocdepth}{2}

% --- Encabezados y pies ---
% A la izquierda, "Proyecto Integrador -- Thoraxis"; a la derecha, el programa
% y el logo de UTEC Posgrado; al pie, el número de página. La portada no lleva
% encabezado (titlepage).
%
% El logo se apoya en la línea base, para no tocar la línea del encabezado, y
% \articleHeaderText sube los dos textos hasta el centro del logo: la mitad
% de la altura del logo menos la mitad de la altura de la "P". Así el texto y
% el logo quedan a la misma altura, los textos de ambos lados siguen
% alineados entre sí y el centrado se mantiene aunque cambie el tamaño del
% logo o de la letra.
\newcommand{\articleHeaderText}[1]{%
    \raisebox{\dimexpr 0.5\articleHeaderLogoHeight - 0.5\fontcharht\font`P\relax}{#1}}
\usepackage{fancyhdr}
\fancypagestyle{article}{%
    \fancyhf{}%
    \fancyhead[L]{\footnotesize\articleHeaderText{\articleDocument{} -- \articleShortTitle}}%
    \fancyhead[R]{\footnotesize\articleHeaderText{\articleProgram}\hspace{0.8em}%
        \includegraphics[height=\articleHeaderLogoHeight]{utec-posgrado}}%
    \fancyfoot[C]{\thepage}%
    \renewcommand{\headrulewidth}{0.4pt}%
}

% --- Figuras y tablas ---
\usepackage{float}          % [H] para fijar una figura en su lugar
\usepackage{array}
\usepackage{tabularx}
\usepackage{booktabs}
\usepackage{multirow}
% Más aire entre filas de las tablas, para que las de varias líneas no se
% junten.
\renewcommand{\arraystretch}{1.25}
\usepackage[
    labelfont=bf,
    labelsep=period,
    font=small,
    justification=centering,
    singlelinecheck=false,
    skip=6pt
]{caption}

% --- Listas ---
\usepackage{enumitem}
\setlist{itemsep=0.2em, topsep=0.4em}

% --- Matemáticas y unidades ---
\usepackage{amsmath}
\usepackage{siunitx}
\sisetup{output-decimal-marker={,}, group-separator={\,}}

% --- Páginas horizontales para figuras anchas ---
% pdflscape gira la página en el visor de PDF, además de en la impresión.
\usepackage{pdflscape}
\usepackage{afterpage}

% --- Control de extensión del resumen ejecutivo ---
% refcount lee los números de página de las etiquetas (ver commands.tex).
\usepackage{refcount}

% --- Bibliografía (IEEE) ---
% Estilo numérico IEEE (biblatex-ieee): las referencias se numeran en el orden
% en que se citan. \parencite{clave} da [n] y \textcite{clave}, «Autor et al. [n]».
\usepackage{csquotes}
\usepackage[style=ieee, backend=biber]{biblatex}
\addbibresource{bibliography/references.bib}
% IEEE usa raya corta (en dash) en los rangos de páginas; el estilo en español
% pone un guion.
\DefineBibliographyExtras{spanish}{\renewcommand*{\bibrangedash}{\textendash}}
% biblatex-ieee solo trae sus textos en inglés: estos son sus equivalentes en
% español (número de artículo, fecha de consulta y dirección web).
\DefineBibliographyStrings{spanish}{%
    artno   = {art\adddot\addnbspace n\adddot\textsuperscript{o}},
    urlseen = {Consultado\addcolon},
    url     = {Disponible en},
}

% --- Enlaces ---
% hyperref va casi al final: redefine comandos de los paquetes anteriores.
\PassOptionsToPackage{hyphens}{url}
\usepackage[
    colorlinks=true,
    linkcolor=black,
    citecolor=black,
    urlcolor=articlelink,
    bookmarksnumbered=true,
    pdfborder={0 0 0}
]{hyperref}
\urlstyle{same}
